\chapter{Boolean Analysis --- Degree One}
%%% generated by the blueprint pipeline from TCSlib.BooleanAnalysis.ThresholdFunctions.DegreeOne
\section{Overview}

This module describes degree-one Fourier weights of subcube and Hamming-ball indicators, exponential and tail estimates for linear forms, and bounds for small sets and biased sign-valued functions. It also states the pi-over-two theorem and an improvement of any linear FKN closeness bound.

\section{Declarations}

\begin{lemma}[First-level weight as a coordinate sum]
\lean{BooleanAnalysis.ThresholdFunctions.weightLevel_one_eq_sum_singleton}
\leanok
\uses{BooleanAnalysis.BooleanFunc,BooleanAnalysis.fourierCoeff,BooleanAnalysis.weightLevel}
\difficulty{4}
Let $n\in\mathbb N$ and let $f:\{0,1\}^n\to\mathbb R$ be a real-valued cube function, with uniform expectation $\mathbb E[h]=2^{-n}\sum_{x\in\{0,1\}^n}h(x)$ and Fourier coefficients
$\widehat f(S)=\mathbb E[f(x)\prod_{i\in S}(-1)^{x_i}]$ for $S\subseteq\{0,\ldots,n-1\}$.
The degree-one Fourier weight $W^1[f]=\sum_{|S|=1}\widehat f(S)^2$ satisfies
\[
W^1[f]=\sum_{i<n}\widehat f(\{i\})^2.
\]
\end{lemma}

\begin{lemma}[Correlation with a linear form]
\lean{BooleanAnalysis.ThresholdFunctions.expect_mul_linear_form}
\leanok
\uses{BooleanAnalysis.BooleanFunc,BooleanAnalysis.boolToSign,BooleanAnalysis.chiS_singleton,BooleanAnalysis.expect,BooleanAnalysis.fourierCoeff,BooleanAnalysis.innerProduct}
\difficulty{3}
Let $n\in\mathbb N$, let $f:\{0,1\}^n\to\mathbb R$ be a real-valued cube function, and let $(b_i)_{i<n}$ be real coefficients.
With uniform expectation $\mathbb E[h]=2^{-n}\sum_{x\in\{0,1\}^n}h(x)$ and singleton Fourier coefficients $\widehat f(\{i\})=\mathbb E[f(x)(-1)^{x_i}]$, one has
\[
\mathbb E\!\left[f(x)\sum_{i<n}b_i(-1)^{x_i}\right]
 =\sum_{i<n}b_i\widehat f(\{i\}).
\]
\end{lemma}

\begin{lemma}[Unit linear form with prescribed correlation]
\lean{BooleanAnalysis.ThresholdFunctions.exists_unit_linear_form}
\leanok
\uses{BooleanAnalysis.BooleanFunc,BooleanAnalysis.ThresholdFunctions.expect_mul_linear_form,BooleanAnalysis.ThresholdFunctions.weightLevel_one_eq_sum_singleton,BooleanAnalysis.boolToSign,BooleanAnalysis.expect,BooleanAnalysis.fourierCoeff,BooleanAnalysis.weightLevel}
\difficulty{4}
Let $n\in\mathbb N$ and let $g:\{0,1\}^n\to\mathbb R$ be a real-valued cube function; write $\mathbb E[h]=2^{-n}\sum_{x\in\{0,1\}^n}h(x)$, $\widehat g(\{i\})=\mathbb E[g(x)(-1)^{x_i}]$, and $W^1[g]=\sum_{i<n}\widehat g(\{i\})^2$ for its degree-one Fourier weight.
If $W^1[g]>0$, there exist real coefficients $(a_i)_{i<n}$ such that
\[
\sum_{i<n}a_i^2=1
\quad\text{and}\quad
\mathbb E\!\left[g(x)\sum_{i<n}a_i(-1)^{x_i}\right]=\sqrt{W^1[g]}.
\]
\end{lemma}

\begin{theorem}[Subcube mean and first-level weight]
\lean{BooleanAnalysis.ThresholdFunctions.subcube_degreeOne_weight}
\leanok
\uses{BooleanAnalysis.BoolCube,BooleanAnalysis.ThresholdFunctions.subcubeIndicator,BooleanAnalysis.boolToSign,BooleanAnalysis.chiS,BooleanAnalysis.chiS_sq_eq_one,BooleanAnalysis.expect,BooleanAnalysis.fourierCoeff,BooleanAnalysis.fourierCoeff_empty,BooleanAnalysis.fourier_coeff_chi,BooleanAnalysis.innerProduct,BooleanAnalysis.uniformWeight,BooleanAnalysis.weightLevel}
\difficulty{5}
Let $n\in\mathbb N$, let $J\subseteq\{0,\ldots,n-1\}$ be nonempty, let $b\in\{0,1\}^n$, and put $k=|J|$ and $q(x)=\mathbf 1_{\{x_i=b_i\ \text{for every }i\in J\}}$ for $x\in\{0,1\}^n$.
Using uniform expectation $\mathbb E[h]=2^{-n}\sum_xh(x)$, singleton Fourier coefficients $\widehat q(\{i\})=\mathbb E[q(x)(-1)^{x_i}]$, and degree-one Fourier weight $W^1[q]=\sum_{i<n}\widehat q(\{i\})^2$, one has
\[
\mathbb E[q]=\frac{1}{2^k}
\quad\text{and}\quad
W^1[q]=\frac{k}{2^{2k}}.
\]
\end{theorem}

\begin{theorem}[Hamming-ball first-level limits]
\lean{BooleanAnalysis.ThresholdFunctions.hammingBall_degreeOne_limit}
\leanok
\uses{BooleanAnalysis.ThresholdFunctions.hammingBallIndicator,BooleanAnalysis.ThresholdFunctions.standardGaussianPDF,BooleanAnalysis.ThresholdFunctions.standardGaussianTail,BooleanAnalysis.expect,BooleanAnalysis.weightLevel}
For $t\in\mathbb R$ and $n\in\mathbb N$, define the $\{0,1\}$-valued function
\[
h_{n,t}(x)=\mathbf 1_{\left\{t\le\frac{\sum_{i<n}(-1)^{x_i}}{\sqrt n}\right\}},
\qquad x\in\{0,1\}^n,
\]
where the quotient at $n=0$ is taken to be $0$.
Write $\mathbb E_n[h]=2^{-n}\sum_{x\in\{0,1\}^n}h(x)$, $\widehat h(\{i\})=\mathbb E_n[h(x)(-1)^{x_i}]$, and $W^1[h]=\sum_{i<n}\widehat h(\{i\})^2$ for the degree-one Fourier weight; let $\varphi(u)=(2\pi)^{-1/2}e^{-u^2/2}$ and $\overline\Phi(t)=\int_t^\infty\varphi(u)\,du$ be the standard Gaussian density and upper tail.
For every real $t$, both limits hold as $n\to\infty$ through the natural numbers:
\[
\mathbb E_n[h_{n,t}]\longrightarrow\overline\Phi(t)
\quad\text{and}\quad
W^1[h_{n,t}]\longrightarrow\varphi(t)^2.
\]
\end{theorem}

\begin{lemma}[Exponential moment of a unit linear form]
\lean{BooleanAnalysis.ThresholdFunctions.rademacher_mgf_le}
\leanok
\uses{BooleanAnalysis.BoolCube,BooleanAnalysis.boolToSign,BooleanAnalysis.expect,BooleanAnalysis.uniformWeight}
\difficulty{4}
Let $n\in\mathbb N$ and let $(a_i)_{i<n}$ be real coefficients with $\sum_{i<n}a_i^2=1$.
For the linear form $L(x)=\sum_{i<n}a_i(-1)^{x_i}$ on $\{0,1\}^n$, with uniform expectation $\mathbb E[h]=2^{-n}\sum_{x\in\{0,1\}^n}h(x)$, every $t\in\mathbb R$ satisfies
\[
\mathbb E[e^{tL(x)}]\le e^{t^2/2}.
\]
\end{lemma}

\begin{theorem}[Truncated first moment of a linear form]
\lean{BooleanAnalysis.ThresholdFunctions.linear_form_tail_expectation}
\leanok
\uses{BooleanAnalysis.BoolCube,BooleanAnalysis.ThresholdFunctions.expect_add,BooleanAnalysis.ThresholdFunctions.expect_const_mul,BooleanAnalysis.ThresholdFunctions.expect_mono,BooleanAnalysis.ThresholdFunctions.rademacher_mgf_le,BooleanAnalysis.boolToSign,BooleanAnalysis.expect}
\difficulty{5}
Let $n\in\mathbb N$, let $(a_i)_{i<n}$ be real coefficients with $\sum_{i<n}a_i^2=1$, and let $s\in\mathbb R$ satisfy $s\ge1$.
For $L(x)=\sum_{i<n}a_i(-1)^{x_i}$ on $\{0,1\}^n$ and uniform expectation $\mathbb E[h]=2^{-n}\sum_{x\in\{0,1\}^n}h(x)$, one has
\[
\mathbb E\!\left[|L(x)|\,\mathbf 1_{\{|L(x)|\ge s\}}\right]
 \le (2s+2)e^{-s^2/2}.
\]
\end{theorem}

\begin{theorem}[Level-one inequality for small sets]
\lean{BooleanAnalysis.ThresholdFunctions.levelOne_inequality}
\leanok
\uses{BooleanAnalysis.BoolCube,BooleanAnalysis.BooleanFunc,BooleanAnalysis.ThresholdFunctions.exists_unit_linear_form,BooleanAnalysis.ThresholdFunctions.expect_add,BooleanAnalysis.ThresholdFunctions.expect_const_mul,BooleanAnalysis.ThresholdFunctions.expect_mono,BooleanAnalysis.ThresholdFunctions.linear_form_tail_expectation,BooleanAnalysis.ThresholdFunctions.weightLevel_nonneg,BooleanAnalysis.boolToSign,BooleanAnalysis.expect,BooleanAnalysis.weightLevel}
\difficulty{5}
For $n\in\mathbb N$ and a real-valued function $h$ on $\{0,1\}^n$, write $\mathbb E[h]=2^{-n}\sum_{x\in\{0,1\}^n}h(x)$, $\widehat h(\{i\})=\mathbb E[h(x)(-1)^{x_i}]$, and $W^1[h]=\sum_{i<n}\widehat h(\{i\})^2$ for its degree-one Fourier weight.
There exists a real constant $C\ge0$ such that, for every $n\in\mathbb N$, every $f:\{0,1\}^n\to\mathbb R$ satisfying $f(x)\in\{0,1\}$ for every $x$, and every real $\alpha$ with $\mathbb E[f]=\alpha$ and $0<\alpha\le\tfrac12$, one has
\[
W^1[f]\le C\alpha^2\log(\alpha^{-1}),
\]
where $\log$ is the natural logarithm.
\end{theorem}

\begin{theorem}[Pi-over-two bound and stability]
\lean{BooleanAnalysis.ThresholdFunctions.pi_over_two_theorem}
\leanok
\uses{BooleanAnalysis.BooleanFunc,BooleanAnalysis.ThresholdFunctions.IsClose,BooleanAnalysis.ThresholdFunctions.degreeOnePart,BooleanAnalysis.ThresholdFunctions.thresholdSign,BooleanAnalysis.fourierCoeff,BooleanAnalysis.isPmOne,BooleanAnalysis.weightLevel}
For $n\in\mathbb N$ and a real-valued function $f$ on $\{0,1\}^n$, use uniform expectation $\mathbb E[h]=2^{-n}\sum_xh(x)$, singleton Fourier coefficients $\widehat f(\{i\})=\mathbb E[f(x)(-1)^{x_i}]$, degree-one Fourier weight $W^1[f]=\sum_{i<n}\widehat f(\{i\})^2$, and degree-one part $L_f(x)=\sum_{i<n}\widehat f(\{i\})(-1)^{x_i}$; define $\tau(r)=1$ for $r\ge0$ and $\tau(r)=-1$ for $r<0$.
There exists a real constant $C\ge0$ such that, for every $n\in\mathbb N$, every sign-valued function $f:\{0,1\}^n\to\{-1,1\}$, and every real $\varepsilon\ge0$ with $|\widehat f(\{i\})|\le\varepsilon$ for every $i<n$, one has
\[
W^1[f]\le\frac2\pi+C\varepsilon
\quad\text{and}\quad
\left(
\frac2\pi-\varepsilon\le W^1[f]
\ \Longrightarrow\
\Pr[f(x)\ne\tau(L_f(x))]\le C\sqrt\varepsilon
\right),
\]
where the probability is uniform on the cube.
\end{theorem}

\begin{definition}[Linear FKN closeness property]
\lean{BooleanAnalysis.ThresholdFunctions.HasFKNClosenessBound}
\leanok
\uses{BooleanAnalysis.BooleanFunc,BooleanAnalysis.ThresholdFunctions.IsClose,BooleanAnalysis.dictator,BooleanAnalysis.isPmOne,BooleanAnalysis.weightLevel}
For $C\in\mathbb R$, define the linear FKN closeness property $\mathcal F(C)$ as follows: for every $n\in\mathbb N$ with $n>0$, every sign-valued function $f:\{0,1\}^n\to\{-1,1\}$, and every real $\delta$ with $0\le\delta\le1$ and $1-\delta\le W^1[f]$, there exist an index $i<n$ and a sign $s\in\{1,-1\}$ such that
\[
\Pr[f(x)\ne s\,d_i(x)]\le C\delta.
\]
Here $d_i(x)=(-1)^{x_i}$ is the $i$th dictator, the probability is uniform on $\{0,1\}^n$, and $W^1[f]=\sum_{i<n}\widehat f(\{i\})^2$ is the degree-one Fourier weight with $\widehat f(\{i\})=2^{-n}\sum_{x\in\{0,1\}^n}f(x)(-1)^{x_i}$.
\end{definition}

\begin{theorem}[First-level weight of a biased function]
\lean{BooleanAnalysis.ThresholdFunctions.biased_degreeOne_weight}
\leanok
\uses{BooleanAnalysis.BoolCube,BooleanAnalysis.BooleanFunc,BooleanAnalysis.ThresholdFunctions.eq_zero_of_expect_eq_zero,BooleanAnalysis.ThresholdFunctions.exists_unit_linear_form,BooleanAnalysis.ThresholdFunctions.expect_add,BooleanAnalysis.ThresholdFunctions.expect_const,BooleanAnalysis.ThresholdFunctions.expect_const_mul,BooleanAnalysis.ThresholdFunctions.expect_mono,BooleanAnalysis.ThresholdFunctions.expect_nonneg,BooleanAnalysis.ThresholdFunctions.mul_thresholdSign_self,BooleanAnalysis.ThresholdFunctions.rademacher_mgf_le,BooleanAnalysis.ThresholdFunctions.thresholdSign,BooleanAnalysis.ThresholdFunctions.thresholdSign_pm_one,BooleanAnalysis.ThresholdFunctions.weightLevel_nonneg,BooleanAnalysis.ThresholdFunctions.weightLevel_one_eq_sum_singleton,BooleanAnalysis.boolToSign,BooleanAnalysis.chiS,BooleanAnalysis.chiS_singleton,BooleanAnalysis.expect,BooleanAnalysis.fourierCoeff,BooleanAnalysis.fourier_coeff_chi,BooleanAnalysis.innerProduct,BooleanAnalysis.isPmOne,BooleanAnalysis.weightLevel}
\difficulty{6}
Let $n\in\mathbb N$ and let $f:\{0,1\}^n\to\{-1,1\}$ be sign-valued, with uniform expectation $\mathbb E[h]=2^{-n}\sum_{x\in\{0,1\}^n}h(x)$, singleton Fourier coefficients $\widehat f(\{i\})=\mathbb E[f(x)(-1)^{x_i}]$, and degree-one Fourier weight $W^1[f]=\sum_{i<n}\widehat f(\{i\})^2$.
For every real $\delta$ satisfying $0\le1-\delta$ and $1-\delta\le|\mathbb E[f]|$, one has
\[
W^1[f]\le4\delta^2\log\!\left(\frac2\delta\right).
\]
The logarithm is natural, and at $\delta=0$ both $2/\delta$ and $\log(2/\delta)$ are interpreted as $0$.
\end{theorem}

\begin{theorem}[Improvement of linear FKN closeness]
\lean{BooleanAnalysis.ThresholdFunctions.fkn_closeness_improvement}
\leanok
\uses{BooleanAnalysis.BoolCube,BooleanAnalysis.BooleanFunc,BooleanAnalysis.ThresholdFunctions.HasFKNClosenessBound,BooleanAnalysis.ThresholdFunctions.IsClose,BooleanAnalysis.ThresholdFunctions.biased_degreeOne_weight,BooleanAnalysis.ThresholdFunctions.disagreementProbability,BooleanAnalysis.ThresholdFunctions.expect_add,BooleanAnalysis.ThresholdFunctions.expect_const,BooleanAnalysis.ThresholdFunctions.expect_const_mul,BooleanAnalysis.ThresholdFunctions.expect_mono,BooleanAnalysis.ThresholdFunctions.expect_nonneg,BooleanAnalysis.ThresholdFunctions.expect_sub,BooleanAnalysis.ThresholdFunctions.weightLevel_one_eq_sum_singleton,BooleanAnalysis.boolToSign,BooleanAnalysis.chiS,BooleanAnalysis.chiS_singleton,BooleanAnalysis.dictator,BooleanAnalysis.expect,BooleanAnalysis.expectationOperator,BooleanAnalysis.flipBit,BooleanAnalysis.flipBit_flipBit,BooleanAnalysis.fourierCoeff,BooleanAnalysis.innerProduct,BooleanAnalysis.isPmOne,BooleanAnalysis.weightLevel}
\difficulty{7}
For any dimension $d\in\mathbb N$ and real-valued function $h$ on $\{0,1\}^d$, set $\widehat h(\{j\})=2^{-d}\sum_{x\in\{0,1\}^d}h(x)(-1)^{x_j}$ and $W^1[h]=\sum_{j<d}\widehat h(\{j\})^2$ for its singleton Fourier coefficients and degree-one Fourier weight, and let every probability on a cube be uniform.
Let $C\in\mathbb R$ satisfy $C\ge1$, and assume that for every $m\in\mathbb N$ with $m>0$, every $g:\{0,1\}^m\to\{-1,1\}$, and every real $\eta$ with $0\le\eta\le1$ and $1-\eta\le W^1[g]$, there exist $j<m$ and $u\in\{1,-1\}$ with $\Pr[g(x)\ne u(-1)^{x_j}]\le C\eta$.
Then for every $n\in\mathbb N$ with $n>0$, every sign-valued $f:\{0,1\}^n\to\{-1,1\}$, and every real $\delta$ with $0\le\delta\le1$ and $1-\delta\le W^1[f]$, there exist $i<n$ and $s\in\{1,-1\}$ such that
\[
\Pr[f(x)\ne s(-1)^{x_i}]
\le\frac\delta4+16C^2\delta^2\max\!\left\{\log\!\left(\frac1{C\delta}\right),1\right\},
\]
where $\log$ is natural and, at $\delta=0$, $1/(C\delta)$ and its logarithm are interpreted as $0$.
\end{theorem}
